% Header: General study / work / project template
%%%%%%%%%%%%%%%%%%%%%%%%%%%%%%%%%%%%%%%%%%%%%%%%%%%%%%%%%%%%%%%%%%%

\documentclass[11pt,a4paper]{scrartcl}

% Layout
\usepackage[margin=2.5cm]{geometry}
\usepackage[onehalfspacing]{setspace}

% General packages
\usepackage{graphicx}
\usepackage{enumitem}
\usepackage{tcolorbox}
\tcbuselibrary{breakable}
\usepackage[breaklinks=true,colorlinks=true,linkcolor=blue,urlcolor=blue,citecolor=blue]{hyperref}
\usepackage{amsmath,amsthm,amssymb,mathtools}
\usepackage{icomma}
\usepackage[T1]{fontenc}
\usepackage[utf8]{inputenc}

% Language: choose one
% \usepackage[ngerman]{babel}
\usepackage[english]{babel}

% Units / tables / floats / references
\usepackage[locale = UK,space-before-unit=true,per-mode = symbol]{siunitx}
\usepackage{booktabs,multirow}
\usepackage{placeins}
\usepackage[natbib,abbreviate=true,doi=false,style=numeric-comp,giveninits=true,sorting=none]{biblatex}
\usepackage{csquotes}

% Optional bibliography
% \addbibresource{references.bib}

% Graphics paths
\graphicspath{{Bilder/} {images/} {figures/} {chapters/}}

% Optional math shortcuts
\newcommand{\R}{\mathbb{R}}
\newcommand{\N}{\mathbb{N}}
\newcommand{\Q}{\mathbb{Q}}
\newcommand{\set}[1]{\left\{ #1 \right\}}
\newcommand{\abs}[1]{\left| #1 \right|}
\newcommand{\norm}[1]{\left\lVert #1 \right\rVert}
\newcommand{\ol}[1]{\overline{#1}}

% Optional units
\DeclareSIUnit{\dBm}{dBm}

% Box styles
\newtcolorbox{calculationbox}[1][]{
    title=Calculation,
    breakable,
    #1
}

\newtcolorbox{solutionbox}[1][]{
    title=Solution,
    breakable,
    #1
}

\newtcolorbox{answerbox}[1][]{
    title=Answer,
    breakable,
    #1
}

\newtcolorbox{notebox}[1][]{
    title=Notes,
    breakable,
    #1
}

\newtcolorbox{sidecalcbox}[1][]{
    title=Side Calculation,
    breakable,
    #1
}

%%%%%%%%%%%%%%%%%%%%%%%%%%%%%%%%%%%%%%%%%%%%%%%%%%%%%%%%%%%%%%%%%%%
\begin{document}

    \title{Physik 4 - Blatt 06}
    \author{Tobias Laser}
    \date{\today}
    \maketitle
    \vfill
    \thispagestyle{empty}

    \pagenumbering{arabic}

    \paragraph{Aufgabe 1: Spin-Bahn-Wechselwirkung}

    Wie in der Vorlesung bereits besprochen wurde, hat die Spin-Bahn-Wechselwirkung in der Kerphysik einen großen Einfluss auf die Aufspaltung von Energieniveaus. Der entsprechende Korrekturterm ist dabei gegeben durch
    \begin{align*}
        V_{ls}(r,l,s) = V_{ls}(r) \cdot (l \cdot s)
    \end{align*}
    Im Folgenden wollen wir dabei die explizite Form von $V_{sl}(r)$ vernachlässigen. Für die Beeinflussung von Zuständen im Schalenmodell ist vor allem der Term $(l \cdot s)$ entscheidend. Analog zur Atomphysik macht es hier Sinn unter dem Einfluss der Spin-Bahn-Wechselwirkung jedem Zustand einen Gesamtdrehimpuls $j = 1 + s$ zuzuordnen.
    \begin{enumerate}
        \item 
            Zeigen Sie mit Hilfe der Definition von $j$
            \begin{align*}
                l \cdot s = \frac{1}{2}(j^2 - l^2 - s^2)
            \end{align*}
            \begin{calculationbox}
                Gegeben ist der Gesamtdrehimpuls
                \begin{calculationbox}
                    $\vec{j} = \vec{k} + \vec{s}$.
                \end{calculationbox}
                Wir wollen daraus einen Ausdruck für $\vec{l} \cdot \vec{s}$ bekommen. Dafür quadrieren wir beide Seiten:
                \begin{calculationbox}
                    $\vec{j}^2 = (\vec{l}^2 + \vec{s}^2)^2$.
                \end{calculationbox}
                Jetzt multiplizieren wir die Klammer aus. 
                \begin{calculationbox}
                    $(\vec{l} + \vec{s})^2 = \vec{l}^2 + 2\vec{l}\cdot \vec{s} + \vec{s}^2$.
                \end{calculationbox}
                Damit gilt:
                \begin{calculationbox}
                    $\vec{j}^2 = \vec{l}^2 + 2 \vec{l} \cdot \vec{s} + \vec{s}^2$.
                \end{calculationbox}
                Stellen nach $\vec{l} \cdot \vec{s}$ um. 
                \begin{calculationbox}
                    $\vec{l} \cdot \vec{s} = \frac{1}{2}(\vec{j}^2 - \vec{l}^2 - \vec{s}^2)$
                \end{calculationbox}
            \end{calculationbox}

        \item
            Betrachten Sie im Folgenden die $1d$-Schale. Welchen Bahndrehimpuls $l$ besitzt diese? In welche Werte von $j$ wird sie durch die Spin-Bahn-Wechselwirkung aufgespalten?
            \begin{calculationbox}
                Wir betrachten die $1d$-Schale.

                Die Bezeichnung besteht aus zwei Teilen:
                \begin{calculationbox}
                    $1d$
                \end{calculationbox}
                Die Zahl $1$ beschreibt hier die radiale Schale. Für den Bahndrehimpuls ist aber der Buchstabe wichtig. 

                Für die Buchstaben gilt:
                \begin{calculationbox}
                    \begin{align*}
                        s \to l = 0\\
                        p \to l = 1\\
                        d \to l = 2\\
                        f \to l = 3\\
                    \end{align*}
                \end{calculationbox}
                Also hat die $1d$-Schale den Bahndrehimpuls
                \begin{calculationbox}
                    $l = 2$
                \end{calculationbox}
                Jetzt kommt der Spin des Nukleons dazu. Protonen und Neutronen haben Spin
                \begin{calculationbox}
                    $s = \frac{1}{2}$
                \end{calculationbox}
                Der Gesamtdrehimpuls insgesamt
                \begin{calculationbox}
                    $\vec{j} = \vec{l} + \vec{s}$
                \end{calculationbox}
                Für die möglichen Werte von $j$ gilt bei Kopplung von $l$ und $s$:
                \begin{calculationbox}
                    $j = l \pm s$
                \end{calculationbox}
                Da $s = \frac{1}{2}$, bekommen wir 
                \begin{calculationbox}
                    $j = l + \frac{1}{2}$ und $j = l - \frac{1}{2}$
                \end{calculationbox}
                Jetzt setzen wir $l = 2$ ein:
                \begin{calculationbox}
                    \begin{align*}
                        j &= 2 + \frac{1}{2} = \frac{5}{2}\\
                        j &= 2 - \frac{1}{2} = \frac{3}{2}\\
                    \end{align*}
                \end{calculationbox}
                Damit spaltet die $1d$ Schale durch die Spin-Bahn-Wechselwirkung in zei Unterschalen auf:
                \begin{calculationbox}
                    $1d_{\frac{5}{2}}$ und $1d_{\frac{3}{2}}$ bzw. $l = 2$, $j = \frac{5}{2}, \frac{3}{2}$ 
                \end{calculationbox}
            \end{calculationbox}

        \item 
            Berechnen Sie den Erwartungswert $\left< l \cdot s\right>$ für beide Unterschalen der $1d$-Schale. Geben Sie anschließend die Energiedifferenz der beiden in Abhängigkeit von $V_{ls}(r)$ an.

            \textit{Hinweis:} $\left< j^2 \right> = j (j + 1)$
            \begin{calculationbox}
                Aus Aufgabe 1a haben wir:
                \begin{calculationbox}
                    $\vec{l} \cdot \vec{s} = \frac{1}{2} (\vec{j}^2 - \vec{l}^2 - \vec{s}^2)$
                \end{calculationbox}
                Für Erwartungswerte verwenden wir:
                \begin{calculationbox}
                    \begin{align*}
                        \left< \vec{j}^2 \right> &= j (j + 1)//
                        \left< \vec{l}^2 \right> &= l (l + 1)//
                        \left< \vec{s}^2 \right> &= s (s + 1)//
                    \end{align*}
                \end{calculationbox}
                Damit
                \begin{calculationbox}
                    $\left< \vec{l} \cdot \vec{s} \right> = \frac{1}{2} [j(j + 1) - l (l + 1) - s (s + 1)]$
                \end{calculationbox}
                Für die $1d$-Schale
                \begin{calculationbox}
                    $l = 2$, $s = \frac{1}{2}$
                \end{calculationbox}
                Also:
                \begin{calculationbox}
                    \begin{align*}
                        l (l + 1) &= 2 \cdot 3 = 6\\
                        s (s + 1) &= \frac{1}{2} \cdot \frac{3}{2} = \frac{3}{4}
                    \end{align*}
                \end{calculationbox}
                \textbf{Erste Unterschale:} $1d_{\frac{5}{2}}$
                Hier ist $j = \frac{5}{2}$
                Also 
                \begin{calculationbox}
                    $j (j + 1) = \frac{5}{2} \cdot \frac{7}{2} = \frac{35}{4}$
                \end{calculationbox}
                und setzen ein
                \begin{calculationbox}
                    \begin{align*}
                        \left< \vec{l} \cdot \vec{s} \right>_{\frac{5}{2}} &= \frac{1}{2} [\frac{35}{4} - 6 - \frac{3}{4}]\\
                        &= \frac{1}{2} [\frac{35}{5} - \frac{24}{4} - \frac{3}{4}]\\
                        &= \frac{1}{2} [\frac{8}{4}]
                        &= 1
                    \end{align*}
                \end{calculationbox}
                \textbf{Zweite Unterschale:} $1d_{\frac{3}{2}}$
                Hier ist $j = \frac{3}{2}$
                Also 
                \begin{calculationbox}
                    $j (j + 1) = \frac{3}{2} \cdot \frac{5}{2} = \frac{15}{4}$
                \end{calculationbox}
                und setzen ein
                \begin{calculationbox}
                    \begin{align*}
                        \left< \vec{l} \cdot \vec{s} \right>_{\frac{3}{2}} &= \frac{1}{2} [\frac{15}{4} - 6 - \frac{3}{4}]\\
                        &= \frac{1}{2} [\frac{15}{5} - \frac{24}{4} - \frac{3}{4}]\\
                        &= \frac{1}{2} [-\frac{8}{4}]
                        &= -\frac{3}{2}
                    \end{align*}
                \end{calculationbox}
                \textbf{Energiedifferenz}
                Die spin-Bahn-Wechselwirkung ist gegeben durch:
                \begin{calculationbox}
                    $V_{ls}(r, \vec{l}, \vec{s}) = V_{ls} \cdot (\vec{l} \cdot \vec{s})$
                \end{calculationbox}
                Daher ist die Energieverschiebung proportional zu 
                \begin{calculationbox}
                    $V_{ls}(r)(\vec{l} \cdot \vec{s})$
                \end{calculationbox}
                Für $1d_{\frac{5}{2}}$
                \begin{calculationbox}
                    $E_{\frac{5}{2}} = V_{ls}(r) \cdot 1 = V_{ls}(r)$
                \end{calculationbox}
                Für $1d_{\frac{3}{2}}$
                \begin{calculationbox}
                    $E_{\frac{3}{2}} = V_{ls}(r) \cdot (-\frac{3}{2}) = -\frac{3}{2} V_{ls}(r)$
                \end{calculationbox}
                Die Energiedifferenz ist:
                \begin{align*}
                    \Delta E &= E_\frac{5}{2} - E_\frac{3}{2}\\
                    &= V_{ls}(r) - (-\frac{3}{2}V_{ls}(r))\\
                    &= \frac{5}{2} v_{ls}(r)
                \end{align*}
            \end{calculationbox}
    \end{enumerate}

    \paragraph{Aufgabe 2: Schalenmodell}

    Wie im Rahmen des Schalenmodells nachvollzogen werden kann, ist ${}^{16}$O ein doppeltmagischer Kern, da er aus sowohl $8$ Protonen und neutronen besteht.
    \begin{enumerate}
        \item 
            Berechnen Sie den Abstand zwischen den Neutronenschalen $1p_\frac{1}{2}$ udn $1d_\frac{5}{2}$ für Kerne im Bereich der Massenzahl $A \approx 16$ aus den totalen Bindungsenergien der Atome ${}^{15}$O ($E_B = 111.96$ MeV), ${}16$O ($E_B = 127.62$ MeV) und ${}^{17}$O ($E_B = 131.76$ MeV).

            \begin{calculationbox}
                Gesucht ist der Abstand zwischen den Neutronenschalen

                ${}^{16}$O hat 8 Protonen und 8 Neutronen. Es ist doppelt magisch, also sind die Schalen bis einschließlich $1p_\frac{1}{2}$ gefüllt.

                Wenn man von ${}^{16}$O ein Neutron entfernt, bekomtm man ${}^{15}$O.

                Das entfernte Neutron stammt aus der äußersten besetzten Neutronenschale $1p_{\frac{1}{2}}$.   
                Die dafür nötige Energie ist die Neutronentrennenergie
                \begin{calculationbox}
                    \begin{align*}
                        S_n({}^{16}\text{O}) &= E_B({}^{16}\text{O} - E_B({}^{15}\text{O}))\\
                        &= 125.62 \text{ MeV} - 111.92 \text{ MeV}\\
                        &= 15.66 \text{ MeV}
                    \end{align*}
                \end{calculationbox}

                ${}^{17}$O: Es ist ${}^{16}$O plus ein zusätzliches Neutron. diese zusätzliche Neutron geht in die nächsthöhere freie Neutronenschale $1d_\frac{5}{2}$.

                Die Bindungsenergie dieses zusätzlichen Neutrons ist:
                \begin{calculationbox}
                    \begin{align*}
                        S_n({}^{17}\text{O}) &= E_B({}^{17}\text{O} - E_B({}^{16}\text{O}))\\
                        &= 131.76 \text{ MeV} - 127.62 \text{ MeV}\\
                        &= 4.14 \text{ MeV}
                    \end{align*}                    
                \end{calculationbox}

                Der Abstand zwischen den Schalen ergibt sich aus der differenz dieser beiden Neutronentrennenergien:
                \begin{calculationbox}
                    \begin{align*}
                        \Delta E &= S_n({}^{16}\text{O}) - S_n({}^{17}\text{O})\\
                        &= 15.66 \text{ MeV} - 4.14 \text{ MeV}\\
                        &= 11.52 \text{ MeV}
                    \end{align*}
                \end{calculationbox}
            \end{calculationbox}

        \item
            Wie interpretieren Sie den Unterschied in der totalen Bindungsenergie von ${}^{17}$O und ${}^{16}$F ($E_B = 128.22$ MeV)?
            \begin{calculationbox}
                \begin{calculationbox}
                    \begin{align*}
                        \Delta E_B &= E_B({}^{17}\text{O} - E_B({}^{17}\text{F}))\\
                        &= 131.76 \text{ MeV} - 128.22 \text{ MeV}\\
                        &= 3.54 \text{ MeV}
                    \end{align*}
                \end{calculationbox}
                Beide Kerne haben dieselbe Massenzahl $A = 17$ Aber 
                \begin{calculationbox}
                    \begin{align*}
                        {}^{17}\text{O}:&&Z = 8, &&N=9\\
                        {}^{17}\text{F}:&&Z = 9, &&N=8
                    \end{align*}
                \end{calculationbox}
                Dem ${}^{17}$O wurde im Vergleich zum doppelt magischen ${}^{16}$O ein Neutron hinzugefügt. Bei ${}^{17}$F wurde statdessen ein Proton hinzugefügt, welches ebenfalls in der nächsten Schale sitzt, aber die zusätzliche elektrische Abstoßung der anderen Protonen verspürt. 

                Durch diese Coulomb-Abstroßung ist der Kern weniger stark gebunden.
            \end{calculationbox}

        \item
            Schätzen Sie anhand von $E_B({}^{17}\text{O}) - E_B({}^{17}\text{F})$ den Radius dieser Kerne ab. Nehmen Sie dabei vereinfacht an, dass es sich um homogene Kugeln mit gleichem Radius handelt und benutzen Sie $E_C = \frac{3}{4}\cdot Z \cdot (Z - 1) \alpha \hbar c / R$ (hier wurde im Vergleich zu bisher die Selbstwechselwirkungsenergie der Protonen abgezogen).
            \begin{calculationbox}
                Für ${}^{17}$O: $Z=8$

                und ${}^{17}$F: $Z=9$

                \begin{calculationbox}
                    \begin{align*}
                        \Delta E_C &= \frac{3}{5}\frac{\alpha\hbar c}{R} [9\cdot 8 - 8 \cdot 7]\\
                        &= \frac{3}{4} \frac{16 \alpha \hbar c}{R}
                        &\stackrel{!}{=} 3.54\text{ MeV}
                    \end{align*}
                \end{calculationbox}
                mit
                \begin{calculationbox}
                    $\alpha \hbar c \approx 3.54$ MeV fm    
                \end{calculationbox}
                und umgeformter Gleichung 
                \begin{calculationbox}
                    $R = \frac{\frac{3}{5} \dot 16 \cdot 1.44\text{ MeV fm}}{3.54\text{ MeV}} \approx 3.9$ fm
                \end{calculationbox}
            \end{calculationbox}
        
        \item
            Vergleichen Sie das Ergebnis mit der Abschätzung $R \approx R_0 \dot A^{\frac{1}{3}}$ (mit $R_0 = 1.2$ fm) und diskutieren Sie die Abweichung.
            \begin{calculationbox}
                \begin{calculationbox}
                    $R \approx R_0 \cdot A^{\frac{1}{3}} = 1.2\text{ fm} \cdot 17^{\frac{1}{3}} = 3.08$ fm
                \end{calculationbox}
                Die über die Coulomb-WW ermittelte $R_{\text{Coulomb}}\approx 3.9$ fm ca. $26\%$ größer.

                \[
                R_{\text{Coulomb}} \approx 3.9\,\mathrm{fm}
                \quad > \quad
                R_{0}A^{1/3} \approx 3.1\,\mathrm{fm}
                \]

                \textbf{Warum größer?}

                \begin{itemize}
                    \item In 2c wurde angenommen:
                    \[
                    \Delta E \text{ stammt nur aus der Coulomb-Energie}
                    \]
                    \item Diese Näherung ist zu einfach $\Rightarrow$ andere Effekte fehlen
                \end{itemize}

                \textbf{Fehlende Beiträge:}

                \begin{itemize}
                    \item Kernkräfte / Schalenstruktur
                    \item Unterschiede zwischen Proton und Neutron
                    \item Annahme eines homogenen Kugelkerns
                \end{itemize}

                \textbf{Fazit:}

                \[
                \boxed{
                \text{Das Coulomb-Modell überschätzt den Radius,}
                \quad
                \text{die Größenordnung stimmt jedoch.}
                }
                \]
            \end{calculationbox}
    \end{enumerate}

\end{document}